\documentclass[11pt]{article}

% --- Packages ---
\usepackage[margin=1in]{geometry}
\usepackage{hyperref}
\usepackage{graphicx}
\usepackage{booktabs}
\usepackage{amsmath}
\usepackage{amssymb}
\usepackage{xcolor}
\usepackage{enumitem}
\usepackage{microtype}

\hypersetup{
  colorlinks=true,
  linkcolor=blue!50!black,
  citecolor=blue!50!black,
  urlcolor=blue!50!black,
}

% --- Metadata ---
\title{Auto-interp Labels Conflate Driver and Thermometer Features\\
       in Sparse Autoencoders:\\
       A Case Study on Pythia-160M}
\author{
  Aung Maw\\
  \texttt{irving46764@gmail.com}\\
  \href{https://github.com/OE-GOD/sae-pythia-160m}{github.com/OE-GOD/sae-pythia-160m}
}
\date{May 2026}

\begin{document}

\maketitle

\begin{abstract}
We train a TopK sparse autoencoder (SAE) on the residual stream of Pythia-160M
at layer 6 and characterize its features through nine analyses. Our headline
finding is methodological: SAE features that receive identical labels from
automated interpretation (auto-interp) can have completely different causal roles.
We demonstrate this concretely on two features both labeled ``Newlines'' by an LLM
labeler. The first (f2255) has a logit weight of 0.494 for newline tokens and
produces 19 newlines per 30 tokens when force-activated. The second (f6767) has a
logit weight of 0.042 and produces zero newlines---a thermometer feature that
detects newline contexts without driving newline output. We argue that
\emph{logit weight analysis}---a single matrix multiply per feature---is a cheap
discriminator between driver and thermometer features and should be standard in
SAE evaluations. The driver/thermometer split sits inside a broader pattern:
stable atomic features (replicate across SAE training runs) and unstable
manifold-partition features (do not replicate, but are locally causally critical)
form orthogonal axes of feature characterization.
\end{abstract}

\section{Introduction}

Sparse autoencoders \cite{bricken2023monosemanticity} extract interpretable features
from transformer activations by training an overcomplete dictionary with sparsity
constraints. The standard interpretation pipeline involves: (1) training the SAE,
(2) finding top-activating examples per feature, and (3) generating a feature
label from those examples, either manually or via an LLM judge
\cite{bills2023language}.

This pipeline conflates two distinct properties of a feature:
\begin{enumerate}[leftmargin=*]
  \item \textbf{Activation correlation:} which inputs cause the feature to fire.
  \item \textbf{Causal effect:} which outputs the feature drives when active.
\end{enumerate}

Auto-interp labels are generated from (1) but typically used as if they describe
(2). When the two diverge---a feature reliably fires on tokens of concept X but
does not causally drive X-related behavior---we call the feature a
\emph{thermometer}. The distinction matters for any application that relies on
features being causally meaningful: monitoring for safety-relevant concepts,
steering model behavior via feature manipulation, or constructing circuit-level
mechanistic explanations.

This work demonstrates the driver/thermometer distinction concretely in a TopK
SAE \cite{gao2024scaling} trained on Pythia-160M and proposes \emph{logit weight
analysis} as a cheap discriminator. We further document that the
driver/thermometer split is concentrated in a specific subspace of features---an
``unstable manifold-partition'' cluster---and contrast its behavior with stable
atomic features that replicate cleanly across SAE training runs.

\section{Setup}

\subsection{Model and SAE}

\begin{tabular}{ll}
\toprule
Base model & Pythia-160M (160M parameters, open) \\
Hook & \texttt{blocks.6.hook\_resid\_post} (middle residual stream, $d_{\text{model}} = 768$) \\
SAE width & $N = 16{,}384$ features (8$\times$ overcomplete) \\
Sparsity & TopK with $k = 64$ active features per token \\
Loss & $L_2$ reconstruction only (no $L_1$; TopK gives hard sparsity) \\
Optimizer & Adam, lr $= 10^{-3}$, batch 4096, 20k steps \\
Data & 1M tokens from \texttt{NeelNanda/pile-10k} \\
Compute & Apple M-series via PyTorch MPS, $\sim$\$50 budget \\
\bottomrule
\end{tabular}

Decoder columns are renormalized to unit length after every optimizer step (to
prevent the optimizer from gaming the sparsity metric by rescaling decoder
weights). The bias $\mathbf{b}_d$ is initialized to the data mean. Encoder is
initialized as the decoder transpose.

We chose TopK over vanilla $L_1$ after empirically observing the well-documented
$L_1$ shrinkage failure mode: a 20$\times$ sweep of $\lambda \in \{0.005, 0.02, 0.1\}$
kept $L_0$ at $\sim$1300, far above any reasonable interpretability target.

\subsection{Reconstruction Quality}

\begin{tabular}{lr}
\toprule
Metric & Value \\
\midrule
Variance explained & 98.43\% \\
$L_0$ (active features per token) & 64.0 (TopK enforced) \\
Dead features & 5.2\% (855 / 16{,}384) \\
$\Delta$CE (information loss) & 0.276 nats/token \\
Loss recovered vs. zero-ablation & 95.4\% \\
\bottomrule
\end{tabular}

The 95.4\% loss-recovery is comparable to small-model results from
\cite{bricken2023monosemanticity} ($\sim$90\%) and \cite{lieberum2024gemma}
(80--95\% across Gemma 2 layers).

\subsection{Auto-interp}

We labeled 26 high-magnitude features using the Moonshot Kimi-K2 API on
top-activating examples. 23 of 26 (88.5\%) were classified as monosemantic with
high or medium confidence.

\section{Experiments and Findings}

We ran nine analyses. Table~\ref{tab:experiments} summarizes them.

\begin{table}[h]
\centering
\caption{The nine experiments and their results.}
\label{tab:experiments}
\small
\begin{tabular}{p{2.5cm}p{4cm}p{6cm}}
\toprule
\textbf{Experiment} & \textbf{Question} & \textbf{Result} \\
\midrule
Width comparison & Right SAE size? & VE saturates at $N{=}16$k; $N{=}64$k has 64\% dead features. \\
CE delta & MLP info preserved? & 95.4\% loss recovered. \\
Auto-interp & Features interpretable? & 88.5\% monosemantic by LLM labeling. \\
Coactivation PMI & Do features fire together? & Unstable: PMI $+0.53$; stable: $-0.16$. \\
Ablation & Per-feature importance? & Stable: $\Delta$CE $\approx 10^{-3}$; unstable: $\Delta$CE $\approx +15$ nats. \\
Cluster geometry & Decoder directions cluster? & Newline cluster cos 0.19; random 0.006 (34$\times$). \\
Splitting (rigorous) & Features split across widths? & Atomic: cos $> 0.99$; manifold: cos 0.35--0.50. \\
Causal intervention & Features causally steer? & 2 of 4 high-conf features pass; 1 fails despite identical label. \\
Logit weights & What predicts causal effect? & Drivers vs thermometers visible; $r = 0.43$ with steering. \\
\bottomrule
\end{tabular}
\end{table}

\subsection{Two Kinds of Features}

Cross-referencing the experiments reveals two distinct populations of features.

\textbf{Stable atomic features} (examples: f5747 file paths, f12117 citation
references, f12697 logical operators):
\begin{itemize}[leftmargin=*]
  \item Cross-seed replication: cos $\geq 0.99$.
  \item Cross-width replication: cos $\geq 0.99$.
  \item Activation rate: 30--50\% of tokens.
  \item Ablation impact: $\Delta$CE $\approx 10^{-3}$--$10^{-2}$ nats/token on active tokens.
  \item Coactivation: independent (mean PMI $= -0.16$).
\end{itemize}

\textbf{Unstable manifold-partition features} (the $\sim$16 features labeled
``newline tokens'' variants):
\begin{itemize}[leftmargin=*]
  \item Cross-seed replication: best match cos $< 0.5$ typically.
  \item Cross-width replication: best 64k match has cos $\approx 0.35$--$0.50$.
  \item Activation rate: 1--2\% of tokens.
  \item Ablation impact: $\Delta$CE $\approx +15$ nats on active tokens.
  \item Coactivation: tightly clustered (mean PMI $= +0.53$ within cluster).
  \item Shared subspace: pairwise cos $= 0.19$, vs. 0.006 random ($34\times$ ratio).
\end{itemize}

\textbf{Stability and importance are orthogonal axes.} Stable features replicate
but carry little marginal information per feature. Unstable features do not
replicate but are locally critical.

\subsection{The Newline Cluster is ``One Shared Atom + N Specializations''}

SVD on the 16 newline decoder columns reveals the underlying structure:

\begin{tabular}{lrr}
\toprule
Singular value rank & Cluster value & Random baseline \\
\midrule
1 & 1.98 (24.4\% variance) & 1.10 (7.6\%) \\
2 & 0.97 & 1.09 \\
3--15 & 0.83--0.96 & 0.89--1.07 \\
16 & 0.82 & 0.89 \\
\bottomrule
\end{tabular}

The first singular value dominates dramatically. The remaining 15 are nearly
uniform---the spectral signature of \emph{one shared common direction plus
$N{-}1$ nearly-independent residual directions}, not a low-dimensional
manifold. The shared direction (PC1) captures 24.4\% of cluster variance and
corresponds to a ``newline-detector'' axis common to all cluster members.

\subsection{Driver vs. Thermometer Within the Cluster}

We tested causal intervention on a sample of high-confidence features via
residual-stream steering: for each feature $i$, we add $\alpha \cdot \mathbf{d}_i$
to the residual stream at layer 6 during generation, where $\mathbf{d}_i$ is the
SAE's decoder column for feature $i$ and $\alpha$ is set to the feature's peak
observed activation. We also compute, for each feature, its \emph{logit weight
for newline tokens}: $\max_{t \in T_{\text{NL}}} (\mathbf{d}_i \cdot W_U)_t$
where $T_{\text{NL}}$ is the set of vocabulary tokens containing newlines and
$W_U$ is the model's unembedding matrix.

\begin{table}[h]
\centering
\caption{Driver/thermometer comparison. Two features both labeled ``Newlines'' by
auto-interp diverge by $12\times$ in logit weight and $\infty$ in causal effect.}
\label{tab:drivertherm}
\begin{tabular}{lrrrr}
\toprule
Feature & PC1 alignment & Logit weight (NL) & Newlines induced \\
\midrule
\textbf{f2255 (driver)} & 0.522 & \textbf{0.494} & \textbf{19} \\
f2757 & 0.469 & 0.489 & 2 \\
f15245 & 0.472 & 0.477 & 0 \\
f14584 & 0.462 & 0.360 & 0 \\
f15230 & 0.527 & 0.237 & 0.5 \\
\textbf{f6767 (thermometer)} & 0.525 & \textbf{0.042} & \textbf{0} \\
\bottomrule
\end{tabular}
\end{table}

Both f2255 and f6767 are labeled ``Newline tokens'' by auto-interp. Both fire on
newline-containing tokens. Both have nearly identical PC1 alignment ($\sim 0.52$).
But forcing f2255 on causes the model to output newlines; forcing f6767 on does
nothing.

The logit weight (column 3 of Table~\ref{tab:drivertherm}) cleanly explains the
behavioral divergence. f6767's decoder direction has essentially zero projection
onto newline tokens in the model's output. It is a \emph{thermometer}: it fires
when newlines appear in the input but does not drive newlines in the output.

Pearson correlation between logit weight and steering effect across the 6 tested
features is 0.43---a notable improvement over PC1 alignment (0.38), though
neither metric fully predicts steering success. The remaining variance is likely
attributable to layer normalization, downstream attention dynamics, and other
nonlinear factors not captured by linear logit-attribution.

\subsection{Random-Direction Control}

To rule out the alternative explanation that any large perturbation to the
residual stream produces drift (not specifically feature-direction perturbations),
we ran the steering protocol with random unit vectors at matched magnitude.
Random-direction steering produced generic garbage (word repetition, weird BPE
fragments) but never specifically the labeled concept (newlines, whitespace).
The drift is direction-specific, validating the causal interpretation.

\section{Methodological Recommendation}

Before claiming an SAE feature represents concept X causally:
\begin{enumerate}[leftmargin=*]
  \item Compute the feature's logit weight for X-related tokens
        ($O(d_{\text{model}} \cdot |\mathcal{V}|)$ per feature, milliseconds).
  \item Confirm it is substantially above zero (and ideally above similarly-labeled
        features that fail the test).
  \item For features used in safety-relevant downstream claims, perform an actual
        steering test to confirm causal effect.
\end{enumerate}

Auto-interpretation alone---based on top-activating examples---will reliably
identify thermometer features as ``X-features.'' Steering experiments are
expensive but logit weight analysis is essentially free and discriminates a
meaningful fraction of cases.

\section{Limitations}

\begin{itemize}[leftmargin=*]
  \item Small causal-intervention sample ($n{=}4$ features tested by steering).
        A robust evaluation needs $\sim$30+ features and a quantitative success
        classifier.
  \item Single model (Pythia-160M), single layer (6). Cross-layer features
        (cf. \cite{lindsey2024crosscoders}) are not captured.
  \item The 0.43 correlation between logit weight and steering effect leaves
        substantial variance unexplained.
  \item Manual judgment of steering success via regex newline-counting; an
        automated classifier (e.g., LLM judge) would be more rigorous.
  \item Auto-interp via Kimi K2 has documented weaknesses: it clustered 10+
        distinct features all as ``newline tokens.'' Discrimination required
        follow-up analysis.
\end{itemize}

\section{Reproducibility}

Full code, configuration, and result JSONs at
\url{https://github.com/OE-GOD/sae-pythia-160m}. License MIT. Compute budget
$\sim$\$50 on Apple M-series via PyTorch MPS. SAE training: $\sim$50 minutes.
Full analysis pipeline: $\sim$3 hours.

\section*{Acknowledgments}

Auto-interp via Moonshot's Kimi K2. Built on the methodological foundations of
\cite{bricken2023monosemanticity,gao2024scaling,lieberum2024gemma,rajamanoharan2024gated}.

\bibliographystyle{plainnat}
\bibliography{references}

\end{document}
